\chapter{膨胀宇宙的几何}

要点在于膨胀的宇宙及其几何。目标是先描述空间膨胀意味着什么，然后再追问是什么驱动了膨胀。牛顿的\(F=ma\)宇宙学稍后才会涉及，只有在那之后，几何与动力学才会被整合进完整的Friedmann--Robertson--Walker框架。

\section{膨胀直线上的共动坐标}

从最简单的可能例子开始：一维空间，而非时空。想象一条无限长的直线，上面等间距地分布着星系。坐标\(x\)此时尚未度量物理距离。它标记星系。可以把\(x\)看作星系的名称，而即使整个排列被拉伸或收缩，这个名称也始终附着于该星系。

设从名为\(x=0\)的星系到名为\(x=1\)的星系的物理距离为\(a\)。那么任意两个星系之间的物理距离为
\begin{equation}
d=a\,\Delta x.
\end{equation}
坐标差\(\Delta x\)只是标签之间的差。因子\(a\)将这一标签差转换为由米尺量出的实际距离。

现在把这条直线想象成一根正在被拉伸的弹性线。标签仍随星系一起，但相邻星系之间的距离随时间变化。同样的关系变为
\begin{equation}
d(t)=a(t)\,\Delta x.
\end{equation}
在这个阶段，均匀性仅仅意味着相邻星系在各处都等间距，尽管这一共同间距可能随时间变化。

\section{哈勃定律与不存在中心}

为求两个固定星系的相对速度，保持这一对星系固定，使\(\Delta x\)不随时间变化。求导得到
\begin{align}
\dot d &= \dot a\,\Delta x \\
       &= \frac{\dot a}{a}\,a\,\Delta x \\
       &= \frac{\dot a}{a}\,d.
\end{align}
若定义
\begin{equation}
H(t)\equiv \frac{\dot a}{a},
\end{equation}
则
\begin{equation}
\dot d = H(t)\,d.
\end{equation}
这就是哈勃定律。

历史上的说法``哈勃常数''具有误导性。该量在空间中是常数，而非在时间中。在给定时刻，它对每个星系都相同，但一般而言，它会随着宇宙演化而变化。把它看作哈勃参数更为恰当。

该定律表明，更遥远的星系退行得更快。它并未挑出一个中心。论证中没有任何东西赋予\(x=0\)或任何其他点以特殊地位。每个观测者都看到同样的规则。因此，在这个模型中，膨胀的宇宙并不是从一个特殊位置发出的爆炸。

红移立即进入。如果两个星系在彼此分离时交换光信号，接收到的光会变红。退行速度越大，红移越大。因此，谱线告诉我们退行速度。距离更难测量，而哈勃早期的估计是粗略的。他最初的图并不是一条完美的直线，而是一组带有很大不确定度的散布点。即便如此，几何原理仍在噪声中存留下来：退行速度与距离成正比。

\begin{classroomqa}
\audiencequestion{如果星系相距足够远，哈勃定律是否会使它们运动得比光速还快？}\lecturetimestamp{00:20:23}
\lecturerresponse{是的，在这种退行意义上确实如此。你所说的正确。我们肯定会回到这一点。这是最重要的事情之一。}
\end{classroomqa}

\section{高维、各向同性与时间依赖性\(a(t)\)}

当加入更多空间维度时，本质内容没有变化。在二维中，想象一个由\((x,y)\)标记的方形星系网格，在三维中则由\((x,y,z)\)标记。这些坐标同样是共动标签。它们随星系一起运动，而不是直接度量物理距离。

为什么\(x\)方向上的间距应当与\(y\)方向上的间距相同？答案是观测性的。如果星系在某一方向上系统性地比另一方向更密集，天空在不同方向上看起来就会不同。在足够大的尺度上，情况并非如此。这种方向上的相同性就是各向同性。

在二维中，两个星系之间的距离为
\begin{equation}
d=a\sqrt{(\Delta x)^2+(\Delta y)^2},
\end{equation}
而在三维中为
\begin{equation}
d=a\sqrt{(\Delta x)^2+(\Delta y)^2+(\Delta z)^2}.
\end{equation}
对于固定的一对星系，坐标差保持不变，因此再次求导得到
\begin{equation}
\dot d=\frac{\dot a}{a}\,d.
\end{equation}
哈勃定律不变。

均匀性意味着从一处到另一处相同。各向同性意味着从一个方向到另一个方向相同。两者在小尺度上都不精确。透过银河系看与从银道面外看是不同的，而且在短距离上，单个星系并非均匀分布。但在对许多星系取平均之后，大尺度宇宙在各处、各方向看起来都近似相同。

\begin{classroomqa}
\audiencequestion{如果速度与距离成正比，而距离又在增大，这是否意味着固定的一对星系必定在彼此加速远离？}\lecturetimestamp{00:40:22}
\lecturerresponse{不。哈勃定律说的是，在某一给定时刻，退行速度与距离成正比。它本身并未说明固定的一对星系的速度随时间增大还是减小。这取决于\(a(t)\)的具体时间依赖关系。若\(a(t)\propto t\)，则相邻星系以恒定速率分离。若\(a(t)\propto t^2\)，则固定的一对星系的退行速率随时间增大。若\(a(t)\propto \log t\)，则它随时间减小。定律\(\dot d=(\dot a/a)d\)作为距离的函数始终成立，但加速与减速是另外的问题。}
\end{classroomqa}

长期以来，宇宙学家预期宇宙在减速，因为引力应当减慢初始的膨胀，正如引力减慢向上抛出的物体一样。这里的要点只是：哈勃定律本身并不决定\(\dot a\)或\(\ddot a\)的行为。

\begin{classroomqa}
\audiencequestion{你如何设定标度因子\(a\)的初始值？}\lecturetimestamp{00:46:18}
\lecturerresponse{它的整体归一化是任意的。物理上有意义的不是\(a\)本身的数值大小，而是它如何变化。如果定义另一个标度因子\(b=3a\)，那么\(\dot b/b=\dot a/a\)。将\(a\)乘以一个常数重新标度，改变的是约定或单位，而不是哈勃参数。}
\end{classroomqa}

\section{从空间度规到时空度规}

现在把\(\Delta x\)替换为\(dx\)。相邻点之间的物理距离为
\begin{equation}
ds=a\,dx,
\end{equation}
所以
\begin{equation}
ds^2=a^2\,dx^2.
\end{equation}
这是一维空间度规。\(dx^2\)的系数为
\begin{equation}
g_{xx}=a^2.
\end{equation}
因此，膨胀的宇宙可以这样描述：度规的空间部分依赖于时间。

要从空间过渡到时空，使用相对论性间隔。在一维空间中，
\begin{equation}
d\tau^2=dt^2-a^2(t)\,dx^2.
\end{equation}
对于共动观测者，\(dx=0\)，所以\(d\tau=dt\)。通常的宇宙时是随星系一起运动的时钟所测得的固有时。

此时黑板上的图很有用，因为它把共动世界线、线元和度规记号放在一起。
\begin{figure}[H]
\centering
\includegraphics[width=0.82\textwidth]{lecture_01_figure_02.png}
\caption{膨胀的一维宇宙中的共动世界线以及相应的时空间隔。}
\end{figure}

在三维空间中，同样的模式给出谨慎的推广\footnote{黑板在00:56:38--00:58:51附近部分被遮挡。这个三维形式来自从一维空间到三维空间的口头推广，以及黑板上可见的对角线模式。}
\begin{equation}
d\tau^2=dt^2-a^2(t)\bigl(dx^2+dy^2+dz^2\bigr),
\end{equation}
其度规为对角形式
\begin{equation}
g_{\mu\nu}=\mathrm{diag}\bigl(1,-a^2,-a^2,-a^2\bigr).
\end{equation}
非对角项为零，
\begin{equation}
g_{tx}=g_{ty}=g_{tz}=0.
\end{equation}

均匀性意味着度规系数可以依赖于时间，但不能依赖于空间位置。各向同性意味着三个空间方向以相同方式进入。

\begin{classroomqa}
\audiencequestion{我们不是应该把\(dx\)看作距离吗？}\lecturetimestamp{01:01:18}
\lecturerresponse{不。\(dx\)不是距离。\(a\,dx\)才是距离。\(dx\)只是一个标签间隔。一个有用的思考方式是：想象在某一时刻，有人穿过星系画出一张坐标网格，然后要求这张网格跟随它们。坐标间隔按构造保持固定，而物理距离发生变化。人们可以把距离的单位放进\(x\)而不是放进\(a\)；这是一种约定。重要的是乘积\(a\,dx\)。}
\end{classroomqa}

\begin{classroomqa}
\audiencequestion{\(dt^2\)的度规系数等于\(1\)。这也只是约定吗？}\lecturetimestamp{01:03:24}
\lecturerresponse{是的。这本身不包含物理。如果该系数只依赖于时间，人们可以改变时间变量，使它变成只是\(dt^2\)。同样，人们可以选择使\(c=1\)的单位。这里不是约定的地方在于：该系数不依赖于位置，尽管它可能依赖于时间。}
\end{classroomqa}

\begin{classroomqa}
\audiencequestion{如果把它代入爱因斯坦场方程，对于能量-动量张量你会得到什么？}\lecturetimestamp{01:05:59}
\lecturerresponse{那是下一类问题，但它属于动力学，而不属于当前的讨论。人们可以把这个度规代入爱因斯坦方程，并问什么样的能量-动量分布会使它成为一个解。这非常重要，但目前的话题是几何，而不是动力学。}
\end{classroomqa}

\section{视界隐藏了什么}

大尺度宇宙只有在作为近似时才被当作均匀且各向同性的。在恒星、行星乃至单个星系的尺度上，它当然不是均匀的。在每一次局部观测中，它也并非严格各向同性。

农田给出了恰当的尺度依赖性。在一种尺度上它看起来规则；在远小得多的尺度上，它显然并非如此；在比整块农田更大的尺度上，它又不再显得均匀。宇宙学也是如此。在对许多星系取平均之后，大尺度分布看起来大致均匀。在足够小的尺度上则并非如此。

宇宙比我们能看见的那部分更大，也比我们永远能看见的那部分更大。在实际能够探测的尺度上，除局部变化之外，它是各向同性且均匀的。超出这一范围，就没有直接的观测保证。

\begin{classroomqa}
\audiencequestion{你说过我们永远不会真正直接查明这一点。这是否意味着可能存在一种间接查明的方式？}\lecturetimestamp{01:10:34}
\lecturerresponse{也许吧。如今大多数想法是，在足够大的尺度上宇宙并非均匀，我自己的赌注也是如此。但这是一个我永远不会赢也永远不会输的赌注，因为出于基本的原因，我们无法比现在看得远太多。这就是视界的存在。}
\end{classroomqa}

\begin{classroomqa}
\audiencequestion{难道不能用引力来寻找你所能看见的范围之外的质量吗？}\lecturetimestamp{01:11:53}
\lecturerresponse{不能。一个遥远的质量可以只是让我们进入自由落体。坐在一个封闭的房间里，我们处于太阳的引力场中并绕它运行，但如果不看外面，就没有简单的方法判断太阳在那里。只有潮汐力会揭示它，而如果源足够远，那些潮汐效应可能小到无法探测。}
\end{classroomqa}

因此，视界是直接知识的终极屏障。他甚至几乎确定地说，宇宙在体积上至少比我们永远能看见的区域大一千倍，并把这一说法的理由留到后面。

\section{有界而无边界}

现在来看闭合几何。无限直线是开放的且无界的。线段是有界的，意思是它只有有限的范围。有界并不意味着有边界，而是指有限的范围。

有墙的立方体不行。它违反均匀性，因为靠近墙的人与远离墙的人不等价；它也违反各向同性，因为即使位于中心的人在不同方向上看到的东西也不同。

圆或球的内部也不行。位于中心的人在每一个方向看到相同的东西，但这种对称性只属于一个特殊的点。离开中心，墙在一个方向上近、在另一个方向上远。于是对于一般的观测者，空间既不是均匀的，也不是各向同性的。

在非常大的尺度上，工作假设是宇宙均匀且各向同性。这就是宇宙学原理。它有用，但并不免受修正。

\begin{classroomqa}
\audiencequestion{处处各向同性是否意味着均匀？}\lecturetimestamp{01:18:57}
\lecturerresponse{是的。如果每一个观测者在每一个方向看到相同的东西，那就意味着均匀。但逆命题不成立。均匀本身并不意味着各向同性。一个宇宙可以是均匀的，却仍然有一个优先方向，例如如果它充满了全都指向同一方向的箭头。}
\end{classroomqa}

正确的有界模型是一个无特征球面的表面。一个被限制在该表面上的二维生物会发现，每一处都局部地与其他每一处相同，每一个方向都与其他每一个方向相同，尽管这个世界是有限的。这就是有界而无边界。

\section{圆上的闭合一维宇宙}

为使代数保持简单，把球的例子约化到一个空间维度：圆上的生命。在一维中，各向同性意味着左与右等价。圆恰好具有这一特征。选取一个点并称之为\(\theta=0\)。其他每一个点都用角坐标\(\theta\)标记，以弧度度量。原点一侧紧邻的点既可以用一个小的负角来描述，也可以用一个略小于\(2\pi\)的角来描述。这种周期性是几何的一部分。

圆上的居民能够接触的是周长，而不是外部的半径或中心。他们可以沿着圆行进直到回到同一点，从而测量绕行一周的总距离。如果在环境平面中画出这个圆，也可以引入一个类似半径的量，但那属于嵌入图像，而不属于他们所测量的内禀几何。

黑板上的图在这里很有用，因为它把角度标记和总周长保留在同一幅图中。
\begin{figure}[H]
\centering
\includegraphics[width=0.82\textwidth]{lecture_01_figure_03.png}
\caption{闭合一维宇宙的角度描述，总周长为\(2\pi a\)。}
\end{figure}

一个方便的约定是\footnote{在01：26：40--01：29：20记号在把整个周长本身当作尺度因子与把它写成\(2\pi a\)之间短暂切换。这里采用后一种约定，因为它是被带入弧长公式的那一种。}
\begin{equation}
L=2\pi a,
\end{equation}
其中\(L\)是总周长。

对于小的角间距\(\Delta\theta\)，物理距离为
\begin{equation}
d=a\,\Delta\theta.
\end{equation}
过渡到无穷小量给出
\begin{equation}
ds=a\,d\theta,
\end{equation}
因此
\begin{equation}
ds^2=a^2\,d\theta^2.
\end{equation}
这就是圆上的局部度规。

如果\(a\)依赖于时间，闭合一维宇宙就会膨胀或收缩。对于固定的一对星系，
\begin{equation}
d(t)=a(t)\,\Delta\theta,
\end{equation}
所以
\begin{equation}
\dot d=\dot a\,\Delta\theta
      =\frac{\dot a}{a}\,d.
\end{equation}
哈勃定律在这里同样成立，至少对于相距不至于远到绕圆另一侧变得有歧义的星系对而言。

时空度规变为
\begin{equation}
d\tau^2=dt^2-a^2(t)\,d\theta^2.
\end{equation}
如果\(a\)是常数，时空图像就是一个圆柱。如果\(a\propto t\)，它就是一个圆锥。如果\(a(t)\)变化得更不规则，图像就变得像花瓶。

把图像倒过来看，星系会越来越靠近。如果圆上有\(360\)个星系，它们在晚期会分布得很开，更早时非常靠近，再早则更近。最终它们不再是分开的星系，再早甚至不再是分开的原子。但空间中仍然没有特殊点。

这就是对该模型中大爆炸发生在何处这一问题最干净的回答。当\(a(t)\to 0\)时，每个间隔\(a(t)\,\Delta\theta\)都趋于零。整个宇宙同时在所有地方收缩。星系并不是从外部空间中的某个特权位置飞走。几何本身在改变尺度。在这个意义上，答案是：处处。

\begin{classroomqa}
\audiencequestion{你说时空是弯曲的。如果那是一个圆柱，它还弯曲吗？}\lecturetimestamp{01:38:04}
\lecturerresponse{不。如果它是一个圆柱，它就不弯曲。那是\(a\)不随时间变化的情形。它看起来弯曲只是因为它被画在更高维空间中。如果把它剪开，它可以不拉伸地铺平。}
\end{classroomqa}

从那里讨论转向圆锥。如果\(a\propto t\)，时空图像就是一个圆锥。在远离其尖端处，圆锥在同样的内禀意义上也是平的：它可以展开成一个平面。尖端则不同。它是一个锥形奇点。

下一步是球面宇宙的几何，然后是决定\(a(t)\)如何增长或收缩的动力学。
